% The reviewed study-design wording is retained where the brief's OLD text is absent.
\section{\src{Research design}}\label{sec:design}
\subsection{\src{Research question and prediction task}}\label{sec:prediction-task}

\src{The general question of this study is whether the information available on the day a conditional lending programme is approved can predict whether the programme will later be interrupted, and whether a flexible model such as a random forest predicts this better than logistic regression or a neural network. We use IMF programmes as the sample for this analysis, because MONA records the scheduled and completed reviews of IMF-supported programmes approved since 2002 (Section~\ref{sec:motivation}). We test whether approval-day information can predict if an IMF programme will be interrupted. The unit of observation is one lending programme. The prediction is made at Board approval, using only data known on that date.}

\src{This is a supervised binary classification task. The output equals 1 when fewer reviews were completed than scheduled at approval. It equals 0 when every scheduled review was completed. The inputs are 13 variables in four groups: deal structure, conditionality, track record and macroeconomic conditions. Section~\ref{sec:variables} lists them.}

\src{Following James et al. (2013, pp. 16–17), we estimate f in Pr(Y = 1 | X) = f(X). The estimate is judged on programmes it has not seen during training. Following Shmueli (2010, p. 291), we treat this as a prediction problem. Coefficients and variable importance are not read as causal effects.}

\src{Performance is measured on a later test set, with interrupted programmes as the positive class. We report accuracy; sensitivity, the share of interrupted programmes the model flags; specificity, the share of completed programmes it correctly clears; and balanced accuracy, the mean of sensitivity and specificity (Brodersen et al., 2010, p. 3121). We also report the area under the receiver operating characteristic curve (AUC), the probability that a model ranks a randomly chosen interrupted programme above a randomly chosen completed one, so 0.500 corresponds to random guessing (Fawcett, 2006, p. 867). Interruptions are the majority in the test set (see Section~\ref{sec:sample}). A rule that always predicts interruption therefore gets high accuracy. Its balanced accuracy is only 0.500. Plain accuracy misleads when classes are unequal (He \& Garcia, 2009, p. 1276). Balanced accuracy removes this advantage. The random forest and the neural network are tuned on balanced accuracy for the same reason; the logistic regression is not tuned.}

\src{The IMF pays its financing in tranches. Each tranche after the first requires the Board to complete a scheduled review (IMF, n.d.). The review checks the quantitative targets and structural conditions agreed at approval. When a review is not completed, the next disbursement is withheld. The country then loses financing the programme had scheduled. An approval-day flag would, if accurate, inform programme design at the IMF and risk assessment by creditors; Section~\ref{sec:policy} discusses whether our results support such a flag.}

\subsection{\src{Study design, data sources and train/test split}}\label{sec:data-sources}

\src{The study is a cross-section of IMF programmes. Each row is one programme, described as it stood at approval. We keep facilities with scheduled reviews, including their predecessors. These are the Stand-By Arrangement (SBA), the Extended Fund Facility (EFF), the Extended Credit Facility (ECF), and the Standby Credit Facility (SCF). The ECF and the SCF are concessional facilities of the IMF’s Poverty Reduction and Growth Trust (PRGT). Facilities without reviews and precautionary arrangements are dropped (see Section~\ref{sec:sample}). Approvals run from 2002 to 2016. Review records in MONA start in 2002. Ending the window in 2016 limits overlap with the COVID-19 period.}

\src{No single public source records programme design and review outcomes together. We therefore built the dataset from IMF files. The IMF Financial Data Query Tool gives the history of lending arrangements (IMF, 2026a); we call this source the commitments file. It provides dates, facility type, and access in percent of quota at approval. MONA covers reviews, conditions, and staff macro estimates (IMF, 2026b). We use four MONA files. Reviews lists scheduled and completed reviews, from which we build the output. Combined lists structural conditions and prior actions. Mecon holds staff estimates of growth, inflation, and the current account. Description holds a cancellation flag, used only to check the output. We also downloaded the April 2026 World Economic Outlook (WEO; IMF, 2026c). We also used the World Bank’s World Development Indicators (WDI; World Bank, 2026). The final macro inputs come from MONA; Section~\ref{sec:variables} explains why.}

\src{The commitments file and MONA share no common identifier. Country names are spelled differently, and approval dates can differ by a few days. We normalise country names and match on the closest approval date within 60 days. Blends, in which two facilities (for example an ECF and an EFF) are approved for the same country on the same day, become one programme, because they share one review schedule. Overlapping or successive programmes of the same country are kept as separate observations (Section~\ref{sec:limitations}). Section~\ref{sec:sample} reports the resulting sample.}

\src{We split the sample by approval year. Training covers 2002 to 2013, and the test set covers 2014 to 2016. The test programmes are never used to fit or tune the models. All inputs are taken as they stood on approval day (see Section~\ref{sec:variables}). Hyperparameters are tuned on the training set with 10-fold cross-validation (CV) and leave-one-out cross-validation (LOOCV). A third scheme, country-grouped 10-fold CV (country CV), holds out whole countries in each fold. Many countries have several programmes, so random folds put one country on both sides. Randomly split folds give over-optimistic error estimates and favour overly complex models when observations are dependent (Roberts et al., 2017, p. 914). We use the country-grouped scores to choose the tuned settings for the random forest and the neural network. A prediction for a programme from a model fitted without that programme is called an out-of-fold prediction.}
